\documentclass[a4paper]{article}     
\usepackage{amsfonts}
\usepackage{amsmath}
\usepackage{amssymb}
\usepackage{graphicx}
\usepackage{color}

\newcommand{\Q}{\mathbb{Q}}
\newcommand{\R}{\mathbb{R}}
\newcommand{\llim}{\lim\limits}
\newcommand{\dint}{\displaystyle\int}

\begin{document}

\noindent \large Auteur: Ilyas Sefiani \\
\noindent \large Studierichting: Informatica\\
\noindent \large Oefeningenreeks 6A nr. 7.4\\

\medskip

\normalsize
We hebben $x_{1}-x_{5}=30$ en voor $x_{3}-x_{4}=-24$\\

$\left \{
\begin{array}{l}
x_{1}-x_{1}q^4=30 \\
x_{1}q^2-x_{1}q^3=-24
\end{array}
\right .$\\

$\left \{
\begin{array}{l}
x_{1}(1-q^4)=30 \\
x_{1}(q^2-q^3)=-24
\end{array}
\right .$\\

$\left \{
\begin{array}{l}
x_{1}=\dfrac{30}{1-q^4} \\
x_{1}(q^2-q^3)=-24
\end{array}
\right .$\\

$\left \{
\begin{array}{l}
x_{1}=\dfrac{30}{1-q^4} \\
\dfrac{30}{1-q^4}(q^2-q^3)=-24
\end{array}
\right .$\\

$\left \{
\begin{array}{l}
x_{1}=\dfrac{30}{1-q^4} \\
\dfrac{30q^2-30q^3}{1-q^4}=-24
\end{array}
\right .$\\


$\left \{
\begin{array}{l}
x_{1}=\dfrac{30}{1-q^4} \\
30q^2-30q^3=-24(1-q^4)
\end{array}
\right .$\\

$\left \{
\begin{array}{l}
x_{1}=\dfrac{30}{1-q^4} \\
30q^2-30q^3=-24+24q^4
\end{array}
\right .$\\

$\left \{
\begin{array}{l}
x_{1}=\dfrac{30}{1-q^4} \\
-24q^4-30q^3+30q^2+24=0
\end{array}
\right .$\\

$\left \{
\begin{array}{l}
x_{1}=\dfrac{30}{1-q^4} \\
-4q^4-5q^3+5q^2+4=0
\end{array}
\right .$\\

$\left \{
\begin{array}{l}
x_{1}=\dfrac{30}{1-q^4} \\
(q-1)(q+2)(-4q^2-q-2)=0
\end{array}
\right .$\\

$q=1$ of $q=-2$, maar $1-(1)^4=0$ dus $q=1$ kan ook niet, alleen $q=-2$!\\

$\left \{
\begin{array}{l}
x_{1}=\dfrac{30}{1-q^4} \\
q=-2
\end{array}
\right .$\\

$\left \{
\begin{array}{l}
x_{1}=-2 \\
q=-2
\end{array}
\right .$\\

De eerste drie termen vinden van deze rij:\\

$x_{1}=-2$, $x_{2}=4$ en $x_{3}=-8$\\


\begin{thebibliography}{999}
\end{thebibliography}
\end{document} 